% year: תשע"ט
% semester: א
% moed: ב
% number: 4
% points: 20
% categories: func-analysis
% source: exams/מבחנים/תשעט/מועד ב תשעט.pdf

\begin{question}
(20 נק') חקרו את הפונקציה $f(x)=\dfrac{x^3}{x^2-4}$ ושרטטו את הגרף שלה.

\underline{צריך לחקור}: תחום הגדרה, רציפות, אסימפטוטות, תחומי עליה וירידה, נקודות קיצון, תחומי קמירות וקעירות.
\end{question}

\begin{hint}
שימו לב שהפונקציה אי-זוגית. חילוק פולינומים נותן $f(x)=x+\frac{4x}{x^2-4}$.
\end{hint}

\begin{solution}
\textbf{תחום הגדרה ורציפות.} $x\ne\pm2$. בתחום $f$ רציפה (פונקציה רציונלית). $f(-x)=-f(x)$, כלומר $f$ אי-זוגית והגרף סימטרי ביחס לראשית. חיתוך עם הצירים: רק $(0,0)$.

\textbf{אסימפטוטות.} אנכיות: ב-$x=\pm2$ המונה $\pm8\ne0$ והמכנה מתאפס, לכן:
\[ \lim_{x\to2^\pm}f(x)=\pm\infty,\qquad \lim_{x\to-2^\pm}f(x)=\pm\infty \]
(לדוגמה ליד $-2$: המונה $\approx-8$, והמכנה $x^2-4$ שלילי עבור $x\to-2^+$ וחיובי עבור $x\to-2^-$). לכן $x=2$ ו-$x=-2$ אסימפטוטות אנכיות. משופעת: לפי חילוק פולינומים
\[ f(x)=x+\frac{4x}{x^2-4},\qquad \frac{4x}{x^2-4}\xrightarrow[x\to\pm\infty]{}0, \]
ולכן $y=x$ אסימפטוטה משופעת בשני הכיוונים.

\textbf{עליה וירידה.}
\[ f'(x)=\frac{3x^2(x^2-4)-x^3\cdot2x}{(x^2-4)^2}=\frac{x^2(x^2-12)}{(x^2-4)^2}. \]
נקודות חשודות: $x=0$, $x=\pm2\sqrt3$. $f'>0$ עבור $|x|>2\sqrt3$ – $f$ עולה ב-$(-\infty,-2\sqrt3)$ וב-$(2\sqrt3,\infty)$. $f'<0$ עבור $|x|<2\sqrt3$, $x\ne0,\pm2$ – $f$ יורדת ב-$(-2\sqrt3,-2)$, ב-$(-2,2)$ (כולל $x=0$, שבה $f'=0$ אך הסימן לא מתחלף) וב-$(2,2\sqrt3)$.

\textbf{קיצון.} $x=-2\sqrt3$: מקסימום מקומי, $f(-2\sqrt3)=\frac{-24\sqrt3}{8}=-3\sqrt3$. $x=2\sqrt3$: מינימום מקומי, $f(2\sqrt3)=3\sqrt3$. ב-$x=0$ אין קיצון.

\textbf{קמירות וקעירות.}
\[ f''(x)=\frac{8x(x^2+12)}{(x^2-4)^3}. \]
מכיוון ש-$x^2+12>0$, סימן $f''$ כסימן $\frac{x}{(x^2-4)^3}$, כלומר כסימן $x(x^2-4)$:
\begin{itemize}
  \item $x<-2$: שלילי – קעורה ($\cap$);
  \item $-2<x<0$: חיובי – קמורה ($\cup$);
  \item $0<x<2$: שלילי – קעורה ($\cap$);
  \item $x>2$: חיובי – קמורה ($\cup$).
\end{itemize}
נקודת פיתול: $(0,0)$ (שם $f''$ מחליפה סימן ו-$f$ מוגדרת). ב-$x=\pm2$ $f$ אינה מוגדרת.

\textbf{סקיצה.} ב-$(-\infty,-2)$: הגרף עולה מתחת לאסימפטוטה $y=x$ עד המקסימום $(-2\sqrt3,-3\sqrt3)\approx(-3.46,-5.2)$ ואז יורד ל-$-\infty$ ליד $x=-2$. ב-$(-2,2)$: יורד מ-$+\infty$ דרך נקודת הפיתול $(0,0)$ (עם משיק אופקי) אל $-\infty$. ב-$(2,\infty)$: יורד מ-$+\infty$ עד המינימום $(2\sqrt3,3\sqrt3)\approx(3.46,5.2)$ ואז עולה ומתקרב מלמעלה לאסימפטוטה $y=x$.
\end{solution}
